% Do not change document class, margins, fonts, etc.
\documentclass[a4paper,oneside,bibliography=totoc]{scrartcl}

% some useful packages (add more as needed)
\usepackage[utf8]{inputenc}
\usepackage{graphicx}
\usepackage{latexsym}
\usepackage{amsmath}
\usepackage{amssymb}
\usepackage{tabularx}
\usepackage{booktabs}
\usepackage{algorithm} % you can modify the algorithm style to your liking
\usepackage{algorithmic}
\usepackage{csquotes}
\renewcommand{\algorithmiccomment}[1]{\hfill\textit{// #1}}
\usepackage[usenames,dvipsnames]{xcolor}
\usepackage[colorlinks,citecolor=Green]{hyperref} % you may change/remove the colors
\usepackage{lipsum} % you do not need this
\usepackage{environ}
\usepackage{tikz}
\usetikzlibrary{positioning,shapes.gates.logic.US,shapes.gates.logic.IEC,automata,arrows, backgrounds}
\usepackage{natbib}
\bibliographystyle{chicagoa}
\setcitestyle{authoryear,round,semicolon,aysep={},yysep={,}} \let\cite\citep

% example enviroments and commands (add more as needed)
\newtheorem{definition}{Definition} \newtheorem{proposition}{Proposition}

\newcommand\todo[1]{\textcolor{red}{\framebox{\footnotesize TODO}\expandafter\ifx\expandafter\relax\detokenize{#1}\relax\else\ #1\fi}}


% DO NOT MODIFY
\newlength\oldparskip
\NewEnviron{answer}[1][8cm]{%
  \sbox{0}{\parbox[t][#1][t]{0.99\linewidth}{%
      \setlength{\oldparskip}{\parskip}%
      \setlength{\parskip}{.5em}%
      \BODY%
      ~%
      \setlength{\parskip}{\oldparskip}%
    }}%

  \begin{tikzpicture}
    % inner sep restores the \fbox padding
    \node[inner sep=\fboxsep, outer sep=0pt] (A) {\usebox{0}};
    % border thickness like \fboxrule
    \draw[line width=\fboxrule] (A.south west) rectangle (A.north east);
  \end{tikzpicture}%
  \vskip\baselineskip
}


\begin{document}

\title{DL26, Assignment 1}
\author{Anh Tu Duong Nguyen (anguyea, 2115931)\\
\todo{Student 2 name (login, matriculation number)}}
\date{\today}

\maketitle

\section*{Task 2a)}
\begin{answer}[12cm]
  % YOUR ANSWER HERE
\end{answer}

\section*{Task 2b)}
\begin{answer}[14cm]
  % YOUR ANSWER HERE
\end{answer}

\section*{Task 2c)}
\begin{answer}[14cm]
  % YOUR ANSWER HERE
\end{answer}

\section*{Task 2d)}
\begin{answer}[19cm]
  % YOUR ANSWER HERE
\end{answer}

\section*{Task 2e)}
\begin{answer}[18cm]
  % YOUR ANSWER HERE
\end{answer}

\section*{Task 3a)}
\begin{answer}
  % YOUR ANSWER HERE
  There are 6 quantities computed using our parameters and inputs and targets in this forward pass.
  \begin{itemize}
    \item $z_{1} = W_{1}x$
    \item $z_{2} = W_{2} + b_{1}$
    \item $z_{3} = \sigma(z_{2})$ where $\sigma(x) = \frac{1}{1 + e^{-x}}$
    \item $z_{4} = W_{2}z_{3}$
    \item $\hat{y} = z_{4} + b_{2}$
    \item $l = (y- \hat{y})^{2}$
  \end{itemize}
\end{answer}

\section*{Task 3b)}
\begin{answer}[19cm]
  % YOUR ANSWER HERE
  There are in total 6 gradients to compute in this backward pass. The gradients are computed using backpropagation which is defined as follows:

  For a given operator $f: \mathbb{R}^I \to \mathbb{R}^O$, the gradient going outside of the operator back to the input is:
  \begin{center}
    $\delta_{in} = J_{f}^T \delta_{out}$ where $J_{f}$ is the Jacobian of $f$.
  \end{center}
  In the forward pass of the compute graph, the following computations are performed:
  \begin{itemize}
    \item matmul
    \item logistic activation is applied
    \item add 
  \end{itemize}
  Even though the compute graph figure shows that we have linear layers. The inputs to this linear layer are a matrix and a vector. According to the lecture, this equals
  matmul.
  The resulting gradient rules are for these operations:
  \begin{itemize}
    \item $\delta_{out}v^T$ or $A^T\delta_{out}$ for the matmul operation
    \item $\sigma(x)(1-\sigma(x))\delta_{out}$ for the logistic activation
    \item $\delta_{out}$ for all inputs in the add operation
  \end{itemize}
  Applying these rules of the compute graph, we get the following gradients:
  \begin{itemize}
    \item $\delta_{l} = 1$
    \item $\delta_{\hat{y}} = -2(y-\hat{y})\delta_{l}$
    \item $\delta_{y} = 2(y-\hat{y})\delta_{l}$
    \item $\delta_{z_{4}} = \delta_{\hat{y}}$
    \item $\delta_{b_{2}} = \delta_{\hat{y}}$
    \item $\delta_{W_{2}} = \delta_{z_{4}}z_{3}^T$
    \item $\delta_{z_{3}} = W_{2}^T\delta_{z_{4}}$
    \item $\delta_{z_{2}} = \sigma(z_{2})(1-\sigma(z_{2}))\delta_{z_{3}}$
    \item $\delta_{b_{1}} = \delta_{z_{2}}$
    \item $\delta_{W_{1}} = \delta_{z_{2}}x^T$
    \item $\delta_{x} = W_{1}^T\delta_{z_{1}}$
  \end{itemize}
\end{answer}

\newpage
\input{ai-declaration.tex}

\end{document}
